\documentclass[11pt]{article}
\usepackage[a4paper,margin=25mm]{geometry}
\usepackage[T1]{fontenc}
\usepackage{lmodern,microtype,amsmath,amssymb,amsthm,mathtools,booktabs,longtable,array,enumitem,xcolor,graphicx}
\usepackage{xurl}
\usepackage[colorlinks=true,linkcolor=blue!45!black,urlcolor=blue!50!black,citecolor=blue!45!black]{hyperref}
\usepackage{fancyhdr}
\pagestyle{fancy}\fancyhf{}\fancyhead[L]{Interacting networks and rare first failure}\fancyhead[R]{Astra research construction}\fancyfoot[C]{\thepage}
\setlength{\headheight}{14pt}
\setlength{\parskip}{5pt}\setlength{\parindent}{0pt}
\setlist{itemsep=3pt,topsep=4pt}
\newtheorem{theorem}{Theorem}[section]
\newtheorem{lemma}[theorem]{Lemma}
\newtheorem{proposition}[theorem]{Proposition}
\newcommand{\TV}{\operatorname{TV}}
\newcommand{\E}{\mathbb E}
\newcommand{\Prb}{\mathbb P}
\newcommand{\poly}{\operatorname{poly}}
\newcommand{\ceil}[1]{\left\lceil #1\right\rceil}
\newcommand{\floor}[1]{\left\lfloor #1\right\rfloor}
\newcommand{\eps}{\varepsilon}
\newcommand{\ind}{\mathbf 1}
\title{A finite first-failure curve for interacting networks\\[7pt]\large Constructing rare-event access from pinned Gibbs sampling}
\author{Astra research construction\\\small Complete derivation, finite implementation, and retained limitations}
\date{5 November 2026}
\hypersetup{pdftitle={A finite first-failure curve for interacting networks},pdfauthor={Astra research construction},pdfsubject={Constructive rare-event access and a finite whole-CDF theorem for declared Gibbs networks}}
\begin{document}
\maketitle
\begin{abstract}
We construct relative stationary disconnection counting and conditional-failure sampling for a declared interacting Gibbs network. The hypotheses are efficient sampling under every coordinate pinning and a fixed bound on the ratio of conditional failure surprisal, rather than independence or a bounded ordinary probability ratio. An explicit Dobrushin heat-bath construction supplies the sampling operation. Weighted contraction supplies the failure witnesses; a heavy-prefix sampling-to-counting reduction evaluates each rare witness without observing it; inverse multiplicity removes their overlaps. All operations use finitely many rational operations and random bits. Composing these operations with the preserved whole-CDF theorem gives one finite certificate for every positive rational deadline under the exact equilibrium-conditioned safe start. Work is polynomial at fixed contrast, with no inverse stationary-failure-probability or mean-lifetime factor. A family with exponentially varying conditional odds and exponentially incorrect independent-edge predictions satisfies the hypotheses. This is a computational consequence for an explicitly declared model. It neither acquires unknown rare physical rates nor proves completeness of a natural-system failure model. Historical firstness and practical superiority are not asserted.
\end{abstract}
\setcounter{tocdepth}{1}
{\small\setlength{\parskip}{0pt}\tableofcontents}
\newpage

\section{The completed endpoint and its status}
The selected question is whether the previously completed finite first-failure curve can be computed for a genuinely interacting invariant law, with the rare stationary access constructed rather than supplied as an oracle. The endpoint includes the entire relative CDF, failure-conditioned states, finite confidence, and finite bit work. It is not a new mean-only estimator or a sequence of separately certified deadline queries.

The new operation below removes the product-law premise for a precise interacting class. The general reversible spectral and Palm identities, head--middle--tail certificate, and changing-invariant-law comparison are preserved accepted machinery \cite{prior}. Their proofs are included unchanged in the accompanying archive. Section~\ref{sec:join} specifies every replacement needed to apply them. The static access proof in Sections~\ref{sec:mix}--\ref{sec:static} is supplied here in full.

Three independent constructive questions informed the choice: finite-observation certification of transport coefficients; fixed-law bounds for chaotic long-time observables; and acquisition of nonlinear coordinates with an unequilibrated passive bath. Their proofs, obstructions, primary comparators, and unresolved operations are preserved separately. None is relabeled as the selected endpoint. In particular, an ideal instantaneous clamp does not become a finite physical acquisition protocol by omission.

\section{Declared model and theorem}
\label{sec:model}
Let $G=(V,E)$ be a known connected undirected graph, with $v=|V|\ge2$ and $m=|E|$. Parallel edges are allowed. A bit $x_e=1$ means that edge $e$ has failed. There are $N\ge m$ binary variables in total. Extra variables, if present, are explicitly declared latent variables, with their weights and dynamics included in the input. Failure $B$ is disconnection of the visible network; $A=B^c$.

The invariant law has full support and an explicit rational factorization
\begin{equation}
 \pi(x)=W(x)/Z,\qquad W(x)=\prod_f\psi_f(x_{S_f}),\qquad \psi_f>0.
 \label{eq:gibbs}
\end{equation}
The input length $L$ includes all factor tables and rational bit lengths. Evaluating a full conditional or $W(x)$ therefore has polynomial bit cost. There is no implicit exponential-size factor table. Fixed factor arity and bounded incidence are convenient sufficient descriptions, not necessary if the supplied tables and conditional bounds already have polynomial size.

For every visible edge, supply certified rational bounds
\begin{equation}
0<q_e^-\le \pi(x_e=1\mid x_{-e})\le q_e^+<1,
\qquad q_e^-\ge(q_e^+)^c,
\label{eq:contrast}
\end{equation}
where $c\ge1$ is fixed. This bounds the ratio of $-\log q_e^-$ to $-\log q_e^+$. The ordinary ratio $q_e^+/q_e^-$ can be exponentially large in $L$. One may replace $c$ by a larger integer, so checking the sufficient condition uses rational powers. Bounds may be supplied with a proof; computing exact extrema from an unrestricted succinct factor graph is not silently assumed easy.

For static sampling, assume a certified Dobrushin influence bound. With $i$ the coordinate updated and $j$ the coordinate changed, let
\[
 a_{ij}\ge\sup_{x,y:\,x_{-j}=y_{-j}}
 \TV\bigl(\pi_i(\cdot\mid x_{-i}),\pi_i(\cdot\mid y_{-i})\bigr),
 \quad a_{ii}=0,\quad \max_i\sum_j a_{ij}\le1-\alpha,
 \quad\alpha>0.
\]
All certified upper bounds are part of the finite input. Alternatively, an explicitly costed sampler uniform over every pinning can replace this condition. Fast unconditioned sampling alone is not that interface.

The operating continuous-time generator is rational, succinct, reversible with respect to this same $\pi$, and has computable ordinary updates, an exit-rate upper bound, and a certified gap $\gamma>0$. An earned sufficient construction is
\begin{equation}
 Q=\sum_{i=1}^N r_i(H_i-I)+Q_{\rm extra},\qquad
 H_if=\E_\pi[f\mid x_{-i}],\quad r_i\ge r_{\min}>0,
 \label{eq:dynamics}
\end{equation}
where every extra local move is reversible with respect to $\pi$. Then $\gamma\ge\alpha r_{\min}$. The heat-bath backbone is an actual part of the model, not a proof-only addition to a different device. An independent unconditional bit-refresh backbone would generally preserve the wrong law.

Choose rational $R\ge2\max_x[-Q(x,x)]$ and put $g=\gamma/R\le1$. For the CDF result assume $0<p=\pi(B)\le1/2$ and target exactly $\nu=\pi(\cdot\mid A)$.

\begin{theorem}[Finite interacting-network reliability]
\label{thm:main}
For the declared model, the following finite randomized operations exist.
\begin{enumerate}
\item For rational $0<\eps,\beta\le1/4$, return $\widehat p$ such that $|\widehat p/p-1|\le\eps$ with probability at least $1-\beta$.
\item For rational $0<\tau\le1/4$, always return a state in $B$, with unconditional output law within $\tau$ TV of $\pi(\cdot\mid B)$. Fresh calls are independent.
\item For rational $0<\eps,\delta\le1$, construct one finite certificate that, with probability at least $1-\delta$, returns an interval containing
\[
 F(t)=\Prb_{\pi(\cdot\mid A)}(\tau_B\le t)
\]
at every positive rational $t$, simultaneously, with endpoint ratio at most $1+\eps$.
\end{enumerate}
At fixed $c$, preprocessing, memory, and random-bit work are polynomial in
\[
 L,\ N,\ m,\ v,\ \alpha^{-1},\ g^{-1},\ \eps^{-1},\ \log(1/\delta)
\]
for the CDF, with the analogous accuracy parameters for the static operations. Query work is polynomial in the certificate size and binary deadline length. There is no factor proportional to $p^{-1}$ or the mean lifetime. Dependence on $g^{-1}$ and $\alpha^{-1}$ is on their values. The cut contribution is $v^{O(c)}$, not jointly polynomial in variable $c$.
\end{theorem}

The theorem is computational access to a known model. It does not infer the factors, topology, reversible dynamics, gap, or failure-rule completeness from a short trace. Physical preparation of $\nu$ and observational acquisition of local rates retain their actual exposure and timing costs.

\section{Finite sampling under every pinning}
\label{sec:mix}
Pin any subset of variables. Conditioning restricts the configurations in the influence suprema and removes coordinates, so the same row bound holds on the remaining variables. This applies even when the pinning has exponentially small prior probability.

\begin{lemma}[Fixed finite pinned sampler]
A sample within TV $\eta$ of any pinned law can be generated in
\[
 O\left(N\alpha^{-1}\log(N/\eta)\right)
\]
local heat-bath updates, with polynomial rational bit work and a deterministic finite number of random bits.
\end{lemma}
\begin{proof}
Use a fixed systematic order of free coordinates, beginning from any pin-compatible state. Couple each Bernoulli update using a common uniform variable. If all coordinate disagreement probabilities before a sweep are at most $d$, an updated coordinate has disagreement probability at most $(1-\alpha)d$. During the sweep no coordinate exceeds $d$; after the sweep all free coordinates have been updated, so their maximum is at most $(1-\alpha)d$. Coupling to a stationary chain and taking a union bound gives TV at most $N(1-\alpha)^S\le Ne^{-\alpha S}$ after $S$ sweeps.

Take $S=\ceil{\alpha^{-1}\ceil{\log_2(4N/\eta)}}$. Evaluate each conditional exactly as a rational, then round it down to a dyadic probability with $b$ bits, where $2^{-b}\le\eta/(4NS)$. Each update changes TV by at most $2^{-b}$; telescoping the kernels costs at most $NS2^{-b}$. Mixing and arithmetic error total at most $\eta/2$, with slack. Each Bernoulli draw uses exactly $b$ fair bits. If there are no free coordinates, return the pinned state directly. Pin constraints are always obeyed.
\end{proof}

The standard random-scan version also gives the required gap \cite{hayes}. Write $H=N^{-1}\sum_iH_i$. Under synchronous coordinate choices, the maximum disagreement probability contracts by $1-\alpha/N$ per step, giving TV at most $N(1-\alpha/N)^n$. Since $H$ is self-adjoint and positive semidefinite, applying this estimate to an eigenfunction and taking $n$th roots gives $\operatorname{gap}(H)\ge\alpha/N$. Thus $\sum_i(H_i-I)$ has gap at least $\alpha$. Dirichlet-form comparison proves $\operatorname{gap}(Q)\ge\alpha r_{\min}$ in \eqref{eq:dynamics}. If total update rate $r$ is divided evenly among coordinates, $r_{\min}=r/N$; the factor $N$ is paid.

\section{Counting a rare pinning without observing it}
\label{sec:prefix}
This is an explicit weighted version of the established sampling-to-counting self-reduction \cite{jvv}. It is not a new replacement for that classical principle.

For a pinning $D$, let $Z_D=\sum_{x\models D}W(x)$. The target is relative precision in $Z_D$, not an additive estimate of the possibly tiny probability of $D$. Constants in $W$ fixed by the pinning are retained; dropping them would corrupt partition ratios.

\begin{lemma}[Heavy-prefix partition estimate]
For $0<u\le1$ and $0<\beta<1$, finitely many pinned samples produce a positive rational $\widehat Z_D$ with
\[
 |\log(\widehat Z_D/Z_D)|\le u
\]
with probability at least $1-\beta$. At most $O(N^3u^{-2}\log(N/\beta))$ pinned draws are needed, each at TV tolerance $u/(16N)$. Output bit length is polynomial in $L,N,\log(1/u),\log\log(1/\beta)$.
\end{lemma}
\begin{proof}
Let $h=u/(8N)$ and
\[
 M=\ceil{2h^{-2}\ceil{\log_2(2N/\beta)}}.
\]
At each free coordinate estimate its conditional marginal $p_i$ from $M$ independent draws under the current prefix, each within $h/2$ TV. Hoeffding's inequality and the TV mean bias give $|\widehat p_i-p_i|\le h$ with conditional probability at least $1-\beta/N$. Choose the empirically heavier bit, whose estimated conditional probability is $\widehat h_i=\max(\widehat p_i,1-\widehat p_i)\ge1/2$, and pin it. On success its true conditional probability satisfies $h_i\ge1/2-h\ge3/8$. Therefore
\[
 |\log\widehat h_i-\log h_i|\le h/(1/2-h)\le(8/3)h.
\]
Adaptive prefixes do not invalidate the conditional probability guarantee; a union bound covers at most $N$ stages. At the completed state $x^*$,
\[
 \pi(x^*\mid D)=\prod_i h_i=W(x^*)/Z_D.
\]
Return $\widehat Z_D=W(x^*)/\prod_i\widehat h_i$. The accumulated log error is at most $(8/3)Nh\le u/3<u$. Every empirical factor is a rational at least $1/2$, so exact products and the known rational $W(x^*)$ have the stated finite bit cost. If $D$ pins all variables, $Z_D=W(x^*)$ is exact.
\end{proof}

For a cut $D$, its event $E_D=\{x_e=1\ \forall e\in D\}$ has probability $p_D=Z_D/Z$. Estimate $Z$ once and $Z_D$ for each retained cut with log tolerance $s/2$ and a joint confidence budget. Then the positive rational weights $\phi_D=\widehat Z_D/\widehat Z$ obey $e^{-s}p_D\le\phi_D\le e^sp_D$. No trial must exhibit $E_D$.

A cheaper special case remains available. If conditional odds oscillate by a fixed factor $K$, then $q_e^+/q_e^-\le K$ and
\[
 \pi(x_e=1\mid D)=q_e^+\E_{\pi(\cdot\mid D)}
 \left[\frac{\pi(x_e=1\mid X_{-e})}{q_e^+}\right],
\]
whose normalized integrand lies in $[1/K,1]$. Multiplying such estimates gives a Rao--Blackwell alternative. The heavy-prefix construction is needed for the stronger class where ordinary conditional ratios can be exponential.

\section{Rational cut geometry and a finite witness list}
\label{sec:cuts}
\subsection{Weights without a logarithm oracle}
For $q=q_e^+$, find the nonnegative integer $s=\floor{\log_2(1/q)}$ by rational comparisons with powers of two. Put $y=2^sq\in(1/2,1]$, $z=(1-y)/(1+y)$, and
\begin{equation}
 w_e=\tfrac23s+2z>0.
 \label{eq:rationalweights}
\end{equation}
For $0\le z\le1/3$, $2z\le2\operatorname{atanh}z\le(9/8)2z$, and $2/3\le\log2\le3/4$. Consequently
\begin{equation}
 w_e\le-\log q_e^+\le\tfrac98w_e\le2w_e.
 \label{eq:logbounds}
\end{equation}
These are exact rationally constructed positive weights, including when $q_e^+$ is extremely close to one. Their bit length is polynomial in $L$.

Set $C=\ceil{2c}$ and compute a minimum weighted cut of value $\lambda>0$ with exact rational arithmetic. Standard min-cut algorithms apply directly; rational weights need not be expanded into repeated edges. The sharper choice $C=\ceil{9c/8}$ is also valid, but the bounds below use $\ceil{2c}$ throughout.

The tower property extends the full-conditional bounds to every partial conditioning. Sequentially exposing the failed bits of any edge set $D$ therefore proves
\begin{equation}
 e^{-Cw(D)}\le\pi(E_D)\le e^{-w(D)},\qquad
 p=\pi(B)\ge e^{-C\lambda}.
 \label{eq:witness}
\end{equation}
No independence or sign of correlation is used.

\subsection{The contraction count, including actual finite randomness}
The following argument is Karger's weighted contraction mechanism \cite{karger}, included to make the witness-list cost explicit.

\begin{lemma}[Integer near-cut enumeration]
For integer $a\ge1$, at most $v^{2a}$ nontrivial unoriented cuts have weight at most $a\lambda$. They can all be listed with confidence $1-\beta$ using a finite number of weighted contractions polynomial in $v,m,v^{2a},\log(1/\beta)$ and the weight bit lengths.
\end{lemma}
\begin{proof}
Contract a weighted-random edge until $2a$ vertices remain, or do no contraction if $v\le2a$. Then choose a uniformly random bipartition by fixing one vertex's side and assigning independent fair bits to the others. An empty side is allowed as a rejected output. For any designated cut of weight at most $a\lambda$, while it survives and $h>2a$ vertices remain, total edge weight is at least $h\lambda/2$. Its probability of being removed in the next contraction is at most $2a/h$. Its survival probability is therefore at least
\[
 \prod_{h=2a+1}^v(1-2a/h)=\binom v{2a}^{-1}.
\]
The final partition matches it with probability $2^{1-2a}$. Since $(2a)!\ge2^{2a-1}$, its output probability is at least $v^{-2a}$. If $v\le2a$, direct uniform bipartition satisfies the same bound. Distinct cut outputs are disjoint events, proving the count.

Put $V_a=v^{2a}$. At every contraction use a dyadic categorical approximation whose TV error is at most $1/(4vV_a)$. Cumulative error is at most $1/(4V_a)$, so a designated cut is still output with probability at least $1/(2V_a)$. Categorical rounding to $b$ bits has TV at most $m2^{-b}$; $b=\ceil{\log_2(4mvV_a)}$ suffices. This uses a fixed number of bits per choice, with no unbounded rejection. Repeat
\[
 R_a=\ceil{2V_a\ceil{\log_2(V_a/\beta)}}
\]
times. A union bound over at most $V_a$ cuts gives failure probability at most $\beta$. Deduplicate outputs and retain only cuts whose exact rational weight is at most $a\lambda$. Always include a computed minimum cut. Thus the list is nonempty even on a failed enumeration event. Contractions, sums, and dyadic comparisons have polynomial bit cost.

When $v\le2a$, a simpler implementation enumerates all $2^{v-1}-1$ nontrivial partitions deterministically and retains the light cuts. This costs at most $V_a$ partitions, because $2^{v-1}\le v^{2a}$. It removes unnecessary random repetition for small graphs while preserving the proved polynomial bound.
\end{proof}

\subsection{Two regimes and the omitted probability}
Let $L_v=\ceil{\log_2(2v)}$. If $\lambda<4L_v$, \eqref{eq:witness} gives the explicit rational floor
\begin{equation}
 p\ge p_0=2^{-8CL_v}.
 \label{eq:p0}
\end{equation}
At fixed $C$, $p_0^{-1}=v^{O(C)}$, so ordinary Gibbs draws and rejection suffice.

Otherwise $\lambda\ge4L_v$. For an accuracy parameter $0<h\le1/4$, choose
\begin{equation}
 B_h=\ceil{\log_2(32/h)},\qquad
 a=\ceil{2C+2+2B_h/\lambda},\qquad V_a=v^{2a}.
 \label{eq:threshold}
\end{equation}
Group cuts heavier than $a\lambda$ into weight intervals $(j\lambda,(j+1)\lambda]$. The integer cut count and \eqref{eq:witness} give
\begin{align*}
 \Prb(\text{some omitted cut fails})
 &\le\sum_{j=a}^\infty v^{2(j+1)}e^{-j\lambda}\\
 &\le2\exp\{2\log(2v)-(\lambda-2\log(2v))a\}.
\end{align*}
Here the geometric-series denominator is at least $1/2$, and $\lambda-2\log(2v)\ge\lambda/2$. Dividing by $p\ge e^{-C\lambda}$ and inserting \eqref{eq:threshold} bounds the ratio by
\[
 8v^2e^{-\lambda}e^{-B_h}\le h/256<h/16.
\]
If $U$ is the union of the retained cut events and enumeration succeeds, then
\begin{equation}
 \pi(B\setminus U)\le(h/16)p,
 \qquad \TV(\pi(\cdot\mid U),\pi(\cdot\mid B))\le h/16.
 \label{eq:tail}
\end{equation}
Finally, the number of retained cuts and the enumeration scale obey
\begin{equation}
 J\le V_a\le64v^{4C+6}/h.
 \label{eq:Jbound}
\end{equation}
Indeed $a<2C+3+2B_h/\lambda$ and $4\log v/\lambda<\log2$, so the remaining factor is at most $2^{B_h}\le64/h$. This makes rarity cost depend on input bit length rather than $p^{-1}$.

\section{Overlap correction, counting, and conditional sampling}
\label{sec:static}
Write the nonempty retained events as $E_1,\ldots,E_J$, let $p_j=\pi(E_j)$, and compute positive rational estimates satisfying
\[
 e^{-s}p_j\le\phi_j\le e^sp_j
\]
on one stated preprocessing event. Let $\widehat Z=\sum_j\phi_j$. Propose a state by choosing $j$ with probability $\phi_j/\widehat Z$ and sampling from the pinned law $\pi(\cdot\mid E_j)$. If these component draws are exact, its density is
\[
 \nu(x)=\frac{\pi(x)}{\widehat Z}A(x),\quad
 A(x)=\sum_j\frac{\phi_j}{p_j}\ind_{E_j}(x).
\]
Let $d(x)=\sum_j\ind_{E_j}(x)$. Every proposal has $1\le d(x)\le J$, and
\begin{equation}
 e^{-s}\le A(x)/d(x)\le e^s.
 \label{eq:density}
\end{equation}
This is the established inverse-multiplicity union construction \cite{karpluby}, now with the weights and component samplers explicitly earned.

\subsection{Relative probability estimator}
The reciprocal observation $Y=1/d(X)$ obeys $1/J\le Y\le1$, and
\begin{equation}
 e^{-s}\pi(U)\le\widehat Z\E_\nu Y\le e^s\pi(U).
 \label{eq:unionmean}
\end{equation}
Set $h=\eps$ in \eqref{eq:threshold}, $s=\eps/16$, and allocate $\beta/4$ each to enumeration and the simultaneous partition estimates. Approximate categorical choice and each pinned component draw within TV $\eps/(256J)$ each. Round $Y$ down to a dyadic mesh at most $\eps/(128J)$. The implemented mean differs from its exact value by at most $\eps/(64J)$, hence relative error at most $\eps/64$. It remains at least $1/(2J)$.

Use
\[
 M_U=\ceil{2048J\eps^{-2}\ceil{\log_2(8/\beta)}}
\]
fresh proposals. Multiplicative bounded-variable concentration gives sample-mean relative error at most $\eps/8$ with failure probability at most $\beta/2$. Return $\widehat p=\widehat Z\overline Y$, clipped above at one if desired. Combining the omitted mass $\eps/16$, log weight error $\eps/16$, implementation error $\eps/64$, and sampling error $\eps/8$ gives $|\widehat p/p-1|\le\eps$ for $\eps\le1/4$. All events have total failure probability at most $\beta$.

In the small-$\lambda$ branch, use unpinned samples at TV $\eps p_0/32$ and at least $256p_0^{-1}\eps^{-2}\ceil{\log_2(2/\beta)}$ trials. Direct failure indicators and the same concentration argument suffice. A zero estimate on a failed confidence event may be replaced by the explicit mass floor below; it does not affect the successful-event statement.

\subsection{An unconditional finite conditional-failure sampler}
For desired TV $\tau$, build a fresh cut list with $h=\tau$, enumeration failure probability $\tau/16$, and joint partition-estimation failure probability $\tau/16$. Set $s=\tau/32$. Under exact mixture proposals, acceptance with probability $1/d(X)$ gives a distribution whose density relative to $\pi(\cdot\mid U)$ lies between $e^{-2s}$ and $e^{2s}$, by \eqref{eq:density}. Its TV error is at most $e^{2s}-1\le\tau/8$.

Use categorical and pinned-draw errors totaling at most $\tau/(128J)$. Round the reciprocal acceptance probability down with error at most $\tau/(128J)$. Conditioning amplifies their joint error by at most $2J$, since the ideal acceptance probability is at least $1/J$. The implemented acceptance probability is pointwise at least $1/J-\tau/(128J)\ge1/(2J)$. Cap attempts at
\[
 M_B=\ceil{2J\ceil{\log_2(16/\tau)}}.
\]
The cap probability is at most $\tau/16$. On cap exhaustion return the known all-failed state. On a bad preprocessing event the output still lies in $B$ and contributes at most its probability, $\tau/8$, to TV. Together with \eqref{eq:tail}, mixture bias, finite proposal error, and the cap, total TV is less than $\tau$.

This is an unconditional TV guarantee, not merely a conditional guarantee given a randomly good table. Rebuilding with fresh randomness for each call makes repeated calls independent. A shared table can improve work, but requires a separate outer confidence event; that optimization is not used in the theorem.

In the small-$\lambda$ branch, draw unpinned states within TV $\tau p_0/32$ and reject until failed. Success probability is at least $p_0/2$; conditional TV error is at most $\tau/16$. Cap at $\ceil{(2/p_0)\ceil{\log_2(16/\tau)}}$ attempts and use the all-failed fallback. The cap adds at most $\tau/16$ TV. Thus this branch also always terminates with a failed state and the promised accuracy.

\subsection{Finite arithmetic and the bill}
Categorical selection over $J$ known rational probabilities can use a uniform $b$-bit grid and rounded cumulative thresholds, with TV at most $J2^{-b}$. For the count estimator, $2^{-b}\le\eps/(256J^2)$ suffices. Reciprocal observations use a common dyadic denominator, so their sums need only $O(b+\log M_U)$ bits. Exact rational partition estimates and their sums have polynomial bit length in $L,N,J$ and the logarithms of trial counts. No exact continuous random variable is an oracle.

The explicit mass floor
\begin{equation}
 \zeta=2^{-N}\prod_f\frac{\min_s\psi_f(s)}{\max_s\psi_f(s)}
 \label{eq:massfloor}
\end{equation}
satisfies $\pi(x)\ge\zeta$ for every state. Every positive conditional event under a feasible pinning also has mass at least $\zeta$. Its reciprocal has polynomial bit length; no loop of length $1/\zeta$ is used.

For a rare-branch construction with requested accuracy $h$, let $V_a$ be \eqref{eq:Jbound}. Enumeration uses $O(V_a\log(V_a/\beta))$ contractions. Estimating all masses uses
\[
 O\bigl((J+1)N^3s^{-2}\log(N(J+1)/\beta)\bigr)
\]
pinned draws. Each such draw costs the proved mixing work in Section~\ref{sec:mix}. Count proposals number $O(J\eps^{-2}\log(1/\beta))$; conditional-output proposals number $O(J\log(1/\tau))$. Each proposal scans the $J$ cut predicates, costing $O(Jm)$ bit tests, in addition to its pinned draw. Stored cut descriptions cost $O(Jm)$ bits, plus polynomial-size rational tables. These conservative bounds can be very large. No near-linear running-time claim is made.

\section{Joining to the finite whole-CDF certificate}
\label{sec:join}
The accepted construction \cite{prior} used a product law only for static counting, conditional-failure draws, safe preparation, a mass floor, and the original refresh-gap proof. Sections~\ref{sec:mix}--\ref{sec:static} replace all five. The spectral and Palm identities themselves use reversibility and a gap, not independence.

Let $P=I+Q/R$, $K=I-P_A$, $u=Rt$, $a_0=\pi(A)$, and $j=\sum_{x\in A,y\in B}\pi(x)P(x,y)$. In the exact safe-conditioned start,
\[
 F(t)=\sum_i\omega_i(1-e^{-\lambda_i u}),\quad
 \omega_i\ge0,\ \sum_i\omega_i=1,
\]
and the preserved facts are
\[
 gpa_0\le j\le p,\qquad gp\le\lambda_1\le p/a_0,\qquad
 \lambda_2\ge g,
\]
\[
 0\le\frac{F(t)-(1-e^{-\lambda_1u})}{1-e^{-\lambda_1u}}
 \le\frac2{gu}.
\]
They concern relative failure probability, not relative accuracy of a vanishing survival tail.

From a conditional failed state, one ordinary $P$ step enters safety with probability $j/p\ge g/2$. Reversing the accepted edge gives a stationary boundary crossing. Forcing one such crossing into a Poisson skeleton of mean $u$ and dividing by the number of crossings gives the bounded Palm identity
\[
 a_0F(t)=pu\,(j/p)\,\E Y_u.
\]
Here $Y_u$ is the indicator that the forced path starts safe, divided by its number of safety-to-failure crossings. Its mean is at least $a_0g/(1+gpu)$. At deadlines where this becomes small, ordinary safe-start simulation has an inverse-polynomial failure probability. A short reverse no-return experiment estimates the slow killed eigenvalue in the rare branch. This is why the full time horizon is not simulated.

\subsection{The exact replacement tolerances}
Use the prior certificate's
\[
 \xi=\eps/40,\quad q_0=g/4,\quad
 \tau_B=\xi gq_0/128,\quad \tau_P=\xi q_0/64,
\]
and a reciprocal-observation dyadic mesh at most $\xi q_0/64$. Supply conditional-failure draws from the unconditional sampler of Section~\ref{sec:static} at tolerance $\tau_B$. The boundary conditioning bias is at most $4\tau_B/g$. The Palm mean bias, including finite Poisson truncation and reciprocal rounding, remains at most $\xi q_0/16$.

For computational safe preparation, draw the unpinned Gibbs law within
\[
 \eta_{\rm eq}=\xi q_0/128
\]
TV and reject failed states. Safe acceptance is at least $1/4$. The accepted law is within $4\eta_{\rm eq}=\xi q_0/32$ of $\pi(\cdot\mid A)$. The direct-simulation mean bias is at most $3\xi q_0/64<\xi q_0/8$. The other preserved bounded-mean biases also remain below $\xi q_0/8$. Thus the prior concentration routine, with ideal mean floor $q_0$, applies unchanged.

For a predetermined total $N_{\rm all}$ of bounded trials, use boundary rejection cap
\[
 C_B=\ceil{(4/g)\ceil{\log_2(8N_{\rm all}/\delta)}},
\]
and the revised safe rejection cap
\[
 C_A=\ceil{4\ceil{\log_2(8N_{\rm all}/\delta)}}.
\]
An outer cap overrun aborts the certificate construction and is charged to its confidence budget. It is not silently substituted into a statistical estimator. Analyze uncapped iid samples, then couple to the finite implementation on the no-abort event. The fresh internal failure samplers already include their own caps in their unconditional TV bounds.

Call the new static count at relative tolerance $\xi/8$ and confidence $\delta/8$, and use \eqref{eq:massfloor} for the prior output rounding. All other finite Poisson tables, rational categorical caps, and confidence allocations remain unchanged. The accepted ledger totals at most $5\delta/8<\delta$, leaving slack. There is no additional confidence allocation for each future query.

\subsection{What one finite object represents}
The certificate stores an initial-slope interval for all $0<u\le\xi$; a finite geometric grid for intermediate times; and either a relative slow-mode interval or a saturation interval for all later times. Concavity of the spectral mixture bounds interpolation between grid points. Finite rational outward evaluation of $1-e^{-\lambda u}$ handles arbitrarily small or large binary-encoded deadlines. The worst interval ratio is $[(1+\xi)/(1-\xi)]^9\le1+\eps$.

The middle horizon is bounded by
\[
 H_{\rm rare}=O\bigl(g^{-1}[\xi^{-1}+\log(1/(g\xi))]\bigr),
\]
\[
 H_{\rm nonrare}=O\bigl(g^{-2}\xi^{-1}\log(1/(g\xi))\log(1/\xi)\bigr).
\]
Let $G_0=O(\xi^{-1}\log(H/\xi))$ be the number of grid points and $M_0=O(g^{-1}\xi^{-2}\log(G_0/\delta))$ the trials per bounded mean. The replacement work is
\begin{align*}
 &T_{\rm count}(\xi/8,\delta/8)
 +O\bigl(M_0(2+G_0C_B)\bigr)T_B(\tau_B)\\
 &\hspace{30mm}+O(M_0G_0C_A)T_{\rm eq}(\eta_{\rm eq}),
\end{align*}
plus the preserved ordinary updates, failure tests, Poisson tables, grid arithmetic, and query evaluator. This is polynomial under the stated parameters, without an inverse lifetime factor. It can be uneconomical; the theorem does not identify these sufficient bounds with optimal costs.

The exact target initial law remains $\nu$. A real preparation merely within absolute TV $\eta$ can perturb $F(t)$ additively by $\eta$, which need not be a relative error as $t\downarrow0$. Computational approximate samples inside bounded identities do not solve that physical initialization problem. The earlier independent-edge all-working point-query extension remains separate; no interacting all-working whole-CDF theorem is claimed here.

\section{Changing invariant law: inherited robustness, explicit Gibbs form}
The preserved continuous-time comparison applies to two killed conductance matrices and safe-set mass matrices satisfying
\[
 e^{-b_L}L_0\preceq L_1\preceq e^{b_L}L_0,
 \qquad e^{-\Delta}M_0\preceq M_1\preceq e^\Delta M_0.
\]
For their respective equilibrium-safe starts it gives
\begin{equation}
 e^{-(b_L+2\Delta)}F_0(t)\le F_1(t)
 \le e^{b_L+2\Delta}F_0(t),\qquad t\ge0.
 \label{eq:uncert}
\end{equation}
This operator result, its sharper shared-backbone common-law form, discrete endpoint qualifications, and physical conditional-update costs are retained in full in the archive.

For an immediate Gibbs consequence, suppose each positive factor entry changes in log magnitude by at most $\eta_f$, and let $E=\sum_f\eta_f$. Normalization gives $|\log(\pi_1(x)/\pi_0(x))|\le2E$; set $\Delta=2E$. If every admitted nonzero transition rate changes by a log factor at most $b$ with common support, each conductance changes by at most $\Delta+b$. Hence $b_L\le\Delta+b$ and
\[
 e^{-(6E+b)}F_0(t)\le F_1(t)\le e^{6E+b}F_0(t).
\]
Combine this with the computed interval $[L(t),U(t)]$ by returning
\[
 [e^{-(6E+b)}L(t),\ \min\{1,e^{6E+b}U(t)\}],
\]
with outward finite arithmetic. Its width can be as large as $(1+\eps)e^{2(6E+b)}$; model uncertainty is not erased by computation. Starting the changed dynamics from the fitted safe law instead of its own safe law adds the separate density-ratio cost $2\Delta$.

This corollary assumes a known common transition support and relative rate bands. Zero-count transitions need positive upper intervals and an appropriate form envelope. The earlier product-law local unknown-support construction is preserved, but it is not automatically a physically acquired Gibbs calibration theorem. No new observational inference of rare factors or omitted external mechanisms is being claimed.

\section{A family where independence is exponentially wrong}
\label{sec:example}
Consider two interacting edge bits with unnormalized weights
\[
 W(0,0)=1,\quad W(1,0)=W(0,1)=t^2,\quad W(1,1)=t^3,
 \qquad t=2^{-b},\ b\ge1.
\]
Their conditional probabilities satisfy
\[
 q^-={t^2\over1+t^2},\qquad q^+={t\over1+t},\qquad
 q^-\ge(q^+)^2.
\]
The conditional odds ratio is $1/t=2^b$, while the Dobrushin row sum is $q^+-q^-<t\le1/2$. Both assumptions of Theorem~\ref{thm:main} have fixed constants despite exponentially different conditional odds.

For two parallel edges, failure means both bits equal one. With $Z=1+2t^2+t^3$, the true static failure probability is $p=t^3/Z$. The independent law with the same one-edge marginals gives $p_{\rm ind}=[t^2(1+t)/Z]^2$, so
\[
 {p\over p_{\rm ind}}={Z\over t(1+t)^2}\asymp2^b.
\]
For $d$ independent such pairs forming $2d$ parallel edges, the error factor is the $d$th power. This is an analytic demonstration of the failure of a product proxy, not the main algorithmic benchmark.

The class is not restricted to independent pairs or bounded interaction degree. On an arbitrary interaction graph $I$ of maximum degree $D\ge1$, let
\[
 W(x)=2^{-2Db\sum_i x_i+b\sum_{ij\in E(I)}x_ix_j},\qquad b\ge1\text{ an integer}.
\]
Every full-conditional failure odds lies between $2^{-2Db}$ and $2^{-Db}$. Valid bounds are
\[
 q^-=\frac1{1+2^{2Db}},\qquad q^+=\frac1{1+2^{Db}},
 \qquad q^-\ge(q^+)^2.
\]
Each neighbor's influence is at most $q^+$, giving row sum at most $D/(1+2^{Db})\le1/2$, uniformly under every pinning. Thus $c=2$ and $\alpha\ge1/2$ suffice for arbitrary declared interaction degree, while conditional odds distortion at a degree-$D$ site is $2^{Db}$. All factors have arity at most two and explicit finite rational bit length; conditional evaluation pays the actual degree. Interaction log strengths grow with $b$. Network topology and interaction topology may be unrelated.

The invariant law can therefore have substantial dependence in the rare configurations that determine failure, even though ordinary safe configurations dominate equilibrium. The algorithm evaluates those configurations through the declared weights and pinning operations rather than replacing their joint probability by marginal products.

\section{Primary alternatives and the exact increment}
\begin{longtable}{>{\raggedright\arraybackslash}p{0.25\textwidth}>{\raggedright\arraybackslash}p{0.32\textwidth}>{\raggedright\arraybackslash}p{0.34\textwidth}}
\toprule
Approach & Established capability & Relation to this construction\\\midrule\endhead
Karger and later network unreliability algorithms \cite{karger,cen} & Relative rare all-terminal static counting for independent edges. The 2025 result is almost linear in the graph size in its common-probability independent model. & Their cut geometry is reused. We supply correlated witness masses and conditional states and join them to the time-CDF theorem. No independent-case speedup is claimed.\\
Sampling--counting self-reduction \cite{jvv} & Approximate counting from appropriate self-reducible sampling. & Heavy-prefix counting is this established machinery made explicit for positive rational Gibbs weights, not a new primitive.\\
Dobrushin/Glauber sampling \cite{hayes} & Rapid mixing under influence conditions; pinning preserves the sufficient bound. & Supplies actual finite sampling and an operating gap. Sampling alone does not estimate exponentially rare disconnection by occupancy.\\
Dependent-reliability splitting \cite{botev,marshall} & Rare-event estimators and conditional sampling for dependent models, including Marshall--Olkin shocks; model-specific efficiency results. & Strong practical and sometimes asymptotic alternatives. The present guarantee is finite worst-case polynomial computation for the declared Gibbs class, with simultaneous relative CDF output.\\
Hedge-connectivity algorithms \cite{hedge} & Combinatorial methods for groups of edges that fail together, with approximation regimes and hardness beyond independent edges. & A different dependence model; these results are not a full Gibbs static oracle or a whole-CDF certificate.\\
Prior finite whole-CDF result \cite{prior} & Finite simultaneous relative first-failure curve with product static access, reversible dynamics, a gap, and exact equilibrium-safe start. & Its time-domain construction is retained. The new operation removes its product static-access premise under explicit correlated sampling and witness bounds.\\\bottomrule
\end{longtable}

Independence is therefore not the essential premise of this rare-event computation. Sufficient conditions are a controlled count of low-cost failure witnesses, probability bounds tied to those costs, efficient conditional sampling under their pinnings, and known finite weights that support sampling-to-counting. For all-terminal failure, the witness enumeration is constructed by contraction. This is a sufficient principle, not a characterization of every tractable dependent reliability model.

No literature search establishes a historical first. The justified claim is the displayed constructive implication and its new scope relative to the preserved product-law endpoint. Nor has a civilization-scale physical capability been demonstrated. The mathematical result is complete under its premises; practical calibration, unmodeled common causes, aging, nonreversible dynamics, and arbitrary initial states remain distinct problems.

\section{Execution and preservation}
\label{sec:evidence}
The static acquisition operation was executed on a nonsymmetric six-vertex network with ten edge variables: five very reliable spanning-tree edges, five moderately unreliable chords, and ten positive pair interactions. One pair interaction has factor $2^{40}$. All factor entries and local conditional tables were represented as exact rationals. Its certified Dobrushin row bound was
\[
 \chi=\frac{1785}{4784}\simeq0.373119,
\]
and the log-contrast condition was checked exactly with $c=3$. Ordinary conditional-odds contrast was about $1.83252\times10^{12}$; this is not a bounded-odds demonstration.

The acquisition code used independently restarted pinned heat-bath chains, heavy-prefix partition estimates, exact local conditional envelopes when sufficiently narrow, and inverse multiplicity. The input to that algorithm was the declared factor model and graph. An isolated reference module enumerated $2^{10}=1024$ states solely to check answers. It did not supply partition functions or conditional samples to the estimator.

\begin{center}
\begin{tabular}{lr}
\toprule
Quantity & Executed value\\\midrule
Exact stationary disconnection probability & $3.72529569\times10^{-28}$\\
Acquired probability estimate & $3.80086379\times10^{-28}$\\
Observed relative estimate error & $2.0285\%$\\
Reported probability interval & $[2.89428757,\,4.75403593]\times10^{-28}$\\
Total pinned draws & $437\,504$\\
Total heat-bath updates & $35\,080\,166$\\
Conditional failed outputs & $8192$\\
Empirical conditional histogram TV distance & $0.0184463$\\
Largest conditional marginal error & $0.0085651$\\
Total observed run time in this environment & $2.99$ seconds\\\bottomrule
\end{tabular}
\end{center}

The reported interval has endpoint ratio approximately $1.6426$. Its local confidence budget is $0.01$ for all partition stages together and $0.01$ for the reciprocal mean, plus separately bounded finite-randomness abort events. Under the mathematical independent-fair-bit model, this supports a confidence of at least $0.98$ minus those abort budgets for this static interval. The executed stream was a reproducible seeded PRNG diagnostic. The observed $2.03\%$ error is not a $2.03\%$ confidence certificate.

Likewise, empirical histogram TV includes finite-output sampling noise and does not certify the output law. Conditional on the acquired partition table's successful confidence event, its proved one-output weight error bound was about $0.145675$, with an additional mixing bound $6.2\times10^{-7}$. The static run reused one table; its setup failure and any joint-output guarantee must be accounted for explicitly. It did not execute the theorem's fresh-preprocessing unconditional-TV sampler at an arbitrarily small target tolerance.

The exact reference also checked every local conditional, the graph identity $d(x)=2^{k(x)-1}-1$ where $k(x)$ is the number of surviving connected components, and the inverse-multiplicity probability identity using rational arithmetic. All 8192 observed probability proposals had multiplicity one, so the run itself did not statistically stress frequent cut overlap; the finite-state identity check included overlapping states. Replacing the correlated law by independent edges with the same marginals gave $1.97938214\times10^{-28}$, underestimating this fixture by a factor about $1.88205$.

Actual random-number counts were $70\,160\,332$ native integer operations consuming $97\,645\,560$ 64-bit PRNG words, plus $32\,768$ Python integer draws consuming $41\,428\,350$ PRNG bits. The native random-cap abort bound was below $7.05\times10^{-14}$; the separate Python bound was below $7.89\times10^{-23}$. Chains used at most 335 heat-bath steps per draw. The native backend checked that every transition denominator fitted in 128 bits and has an arbitrary-precision Python fallback. Local table preparation pays exponential neighborhood enumeration in this small implementation; the general theorem instead permits direct factor evaluation and supplied certified influence bounds.

All 31 cuts were explicitly enumerated for this small fixture, using the direct small-graph alternative justified in Section~\ref{sec:cuts}. At $h=0.1$, the proved cutoff has $v\le2a$ and retains every cut. The general weighted-contraction branch was not implemented or executed by this demonstration. No dynamics or whole-CDF simulation was run. A deliberately conservative schedule calculated for simultaneous 10-percent partition estimates on these 32 partition queries would require up to $26\,214\,400\,000$ pinned draws and $6\,291\,456\,000\,000$ updates; it was not executed, and it is not the full CDF bill. The run used its actual reported intervals and local envelopes instead.

Exact enumeration was faster on this tiny reference problem: its reference and identity checks took about $0.10$ seconds. The experiment establishes execution of the new static representation on a genuinely correlated rare model. The general polynomial consequence comes from the proof, not a speedup over this small exact solver. Scripts, seeds, rational results, source, and reproduction instructions are preserved in the static-execution directory.


The archive preserves every byte of the previous 86-file construction package, including its nested 47-file deferred-calibration package. A SHA-256 manifest records the source and copied bytes. This retains the complete passive classical-output Hamiltonian calibration, its finite receiver, inverse, connected-set compiler, and exposure bound; the finite whole-CDF theorem, changing-law uncertainty, and separate all-working point-query extension; and the earlier all-strength quantum acquisition laws, matched parity/GHZ law, finite-slew and calibration constructions, output-only recovery, bounded-neighborhood result, and thermal-effusion construction at their original scopes. Their source access, budgets, degree, initialization, shell geometry, force/clearing interface, exposure, and implementation limits are not weakened by this new theorem.

An exact model evaluation is distinct from a physical measurement. A factor or local rate of order $2^{-b}$ may be represented and used in polynomial bit work while estimating that rate from physical dwell can still require order $2^b$ exposure. Static Gibbs pinning is a computation on the declared model, not a promise that a physical apparatus can be clamped into every rare configuration at no cost.

\begin{thebibliography}{20}
\bibitem{prior} Accepted Astra construction, \emph{Passive Hamiltonian calibration and finite whole-CDF reliability}, fixed public revision \texttt{76f6a7ce78867b927680ba3b9cd5eb407ce61ccd}. \url{https://github.com/harjassand/astra-research-knowledge/tree/76f6a7ce78867b927680ba3b9cd5eb407ce61ccd/research/returns/2026-11-05-p3-passive-calibration}. Governing original proofs and programs are preserved in the accompanying archive.
\bibitem{karger} D. R. Karger, \emph{A randomized fully polynomial time approximation scheme for the all-terminal network reliability problem}, corrected SIAM Review version, 43 (2001), 499--522. \url{https://people.csail.mit.edu/karger/Papers/reliability-sirev.pdf}.
\bibitem{jvv} M. R. Jerrum, L. G. Valiant, V. V. Vazirani, \emph{Random generation of combinatorial structures from a uniform distribution}, Theoretical Computer Science 43 (1986), 169--188. \url{https://doi.org/10.1016/0304-3975(86)90174-X}.
\bibitem{hayes} T. P. Hayes, \emph{A simple condition implying rapid mixing of single-site dynamics on spin systems}, FOCS 2006. \url{https://www.cs.unm.edu/~hayes/papers/EigenvalueDobrushin/EigenvalueDobrushin-FOCS06.pdf}.
\bibitem{karpluby} R. M. Karp, M. G. Luby, \emph{A new Monte-Carlo method for estimating the failure probability of an n-component system}, UCB/CSD-83-117 (1983). \url{https://www2.eecs.berkeley.edu/Pubs/TechRpts/1983/6351.html}.
\bibitem{cen} R. Cen, J. Li, D. Panigrahi, \emph{Network unreliability in almost-linear time}, STOC 2025. \url{https://arxiv.org/abs/2503.23526}. DOI: 10.1145/3717823.3718140.
\bibitem{botev} Z. Botev, P. L'Ecuyer, G. Rubino, R. Simard, B. Tuffin, \emph{Static network reliability estimation via generalized splitting}, INFORMS Journal on Computing 25 (2013), 56--71. \url{https://doi.org/10.1287/ijoc.1110.0493}.
\bibitem{marshall} Z. I. Botev, P. L'Ecuyer, R. Simard, B. Tuffin, \emph{Static network reliability estimation under the Marshall--Olkin copula}, GERAD G-2015-12 (2015), 28 pages. \url{https://www.gerad.ca/en/papers/G-2015-12}.
\bibitem{hedge} M. Ghaffari, D. R. Karger, D. Panigrahi, \emph{Random contractions and sampling for hypergraph and hedge connectivity}, SODA 2017. \url{https://groups.csail.mit.edu/tds/papers/Ghaffari/soda17-hedge.pdf}.
\end{thebibliography}
\end{document}
